\section{Tanda Tangan Digital}

\begin{subcpmk}
  \item Mahasiswa mampu menganalisis mekanisme tanda tangan digital, membuktikan properti \textit{non-repudiation}, mengevaluasi peran fungsi hash dalam efisiensi tanda tangan, serta memahami integrasi tanda tangan digital dalam ekosistem PKI.
\end{subcpmk}

\noindent\textit{Rujukan:} Stallings \cite{stallings2017}; HAC Bab~8 \cite{hac_online}; NIST SP 800-107.

\subsection{Konsep dan Tujuan Tanda Tangan Digital}

Tanda tangan digital (\textit{Digital Signature}) adalah mekanisme kriptografi yang digunakan untuk mengikat identitas pengirim dengan pesan yang dikirimkan. Berbeda dengan MAC yang hanya memberikan autentikasi antara dua pihak yang berbagi kunci, tanda tangan digital menggunakan \textbf{kriptografi asimetris} untuk memberikan jaminan keamanan yang lebih kuat.

Tanda tangan digital memberikan tiga jaminan keamanan utama:
\begin{enumerate}
    \item \textbf{Integritas (\textit{Integrity})}: Menjamin bahwa pesan tidak mengalami perubahan sejak ditandatangani.
    \item \textbf{Autentikasi (\textit{Authenticity})}: Menjamin bahwa pesan benar-benar dikirim oleh pemilik kunci privat yang bersangkutan.
    \item \textbf{Nir-Sangkalan (\textit{Non-repudiation})}: Pemilik kunci privat tidak dapat menyangkal telah mengirimkan pesan tersebut, karena hanya ia yang memiliki kunci privat untuk membangkitkan tanda tangan yang valid.
\end{enumerate}

\subsection{Mekanisme Kerja Tanda Tangan Digital}

Proses tanda tangan digital menggabungkan fungsi hash dan enkripsi asimetris untuk efisiensi dan keamanan. Alasan utama penggunaan hash adalah karena proses enkripsi asimetris sangat lambat jika diterapkan pada pesan besar.

\subsubsection{Proses Penandatanganan (\textit{Signing})}
\begin{enumerate}
    \item \textbf{Hashing}: Pengirim menghitung nilai hash dari pesan $M$: $h = H(M)$.
    \item \textbf{Signing}: Pengirim mengenkripsi hash tersebut menggunakan \textbf{kunci privatnya} $K_{priv}$: $S = E_{K_{priv}}(h)$.
    \item \textbf{Transmisi}: Pengirim mengirimkan pesan $M$ dan tanda tangan $S$.
\end{enumerate}

\subsubsection{Proses Verifikasi}
\begin{enumerate}
    \item \textbf{Hashing}: Penerima menghitung hash dari pesan $M$ yang diterima: $h' = H(M)$.
    \item \textbf{Verification}: Penerima mendekripsi tanda tangan $S$ menggunakan \textbf{kunci publik} pengirim $K_{pub}$: $h'' = D_{K_{pub}}(S)$.
    \item \textbf{Comparison}: Jika $h' = h''$, maka tanda tangan dinyatakan valid.
\end{enumerate}

\begin{figure}[htbp]
  \centering
  \includegraphics[width=0.8\textwidth]{bab-08/digital-signature-flow}
  \caption{Alur kerja tanda tangan digital: proses hashing, penandatanganan dengan kunci privat, dan verifikasi dengan kunci publik.}
  \label{fig:sig-flow}
\end{figure}

\subsection{Analisis Keamanan: Mengapa Harus Hash Terlebih Dahulu?}

Jika pengirim menandatangani pesan $M$ secara langsung tanpa hashing ($S = E_{K_{priv}}(M)$), maka akan muncul beberapa risiko serius:
\begin{enumerate}
    \item \textbf{In-efficiency}: Enkripsi asimetris sangat lambat untuk pesan besar.
    \item \textbf{Existential Forgery}: Penyerang mungkin dapat memanipulasi cipherteks untuk menghasilkan tanda tangan yang tampak valid untuk pesan yang tidak bermakna.
    \item \textbf{Structure Attack}: Pada RSA, tanda tangan langsung rentan terhadap serangan homomorfik, di mana penyerang dapat menciptakan tanda tangan valid untuk pesan baru dengan mengalikan tanda tangan dari dua pesan yang sudah diketahui.
\end{enumerate}
Dengan menggunakan hash, kita memastikan bahwa yang ditandatangani adalah representasi unik dan kompak dari pesan, sehingga menghilangkan risiko-risiko di atas.

\subsection{Tanda Tangan Digital dalam Ekosistem PKI}

Agar tanda tangan digital dapat dipercaya, penerima harus yakin bahwa kunci publik yang digunakan untuk verifikasi benar-benar milik pengirim. Masalah kepercayaan ini diselesaikan melalui \textbf{Public Key Infrastructure} (PKI).

Sertifikat digital yang diterbitkan oleh \textit{Certificate Authority} (CA) berfungsi sebagai bukti kepemilikan kunci publik. Saat menerima pesan bertanda tangan, penerima melakukan langkah berikut:
\begin{enumerate}
    \item Memeriksa sertifikat digital pengirim.
    \item Memverifikasi tanda tangan CA pada sertifikat tersebut menggunakan kunci publik CA (yang sudah terpercaya/terinstal di sistem).
    \item Setelah sertifikat terverifikasi, gunakan kunci publik pengirim yang ada di dalam sertifikat untuk memverifikasi tanda tangan pesan.
\end{enumerate}

\begin{aktivitas}
  \item \textbf{Audit Sertifikat SSL}: Buka browser, kunjkan situs web dengan HTTPS, dan bedah sertifikat digitalnya. Identifikasi siapa CA yang menandatangani sertifikat tersebut, algoritma hash apa yang digunakan (misal: SHA-256), dan algoritma asimetris apa yang digunakan (misal: RSA atau ECDSA).
\end{aktivitas}

\begin{latihan}
  \item Bedakan secara mendalam antara MAC dan Tanda Tangan Digital dalam hal properti \textit{non-repudiation}! Mengapa MAC tidak bisa memberikan jaminan nir-sangkalan?
  \item Misalkan seorang penyerang menemukan kolisi pada fungsi hash yang digunakan untuk tanda tangan digital. Bagaimana penyerang dapat menggunakan kolisi ini untuk memalsukan dokumen yang sah?
  \item Jelaskan risiko keamanan yang terjadi jika kunci privat pengirim bocor, namun kunci publiknya tetap digunakan oleh penerima untuk memverifikasi pesan lama!
  \item Buktikan bahwa tanda tangan digital menggunakan RSA adalah bentuk invers dari proses enkripsi RSA biasa!
\end{latihan}
